We thank the Area Editor, the Associate Editor, and the referees for the thoughtful, technically detailed, and constructive reports. The feedback was especially valuable in identifying the issues that mattered most for an \emph{Operations Research} revision: the open-system objective, the GPU KV-cache memory mechanism, the role of endogenous memory growth, the proof logic for the threshold policies, and the interpretation of latency evidence. We have therefore revised the manuscript substantially, not only by correcting local notation and exposition issues, but also by reorganizing the model, theory, and experiments around these points. Before responding point by point, we first describe the major changes in the revised manuscript.

\begin{itemize}
    \item \textbf{Reframed the objective with exogenous arrivals.} We revised Sections 2--4 to state the objective under exogenous arrivals more carefully. With fixed external arrivals, a policy in the stable region can handle the offered load. The operational question is whether scheduling can sustain larger arrival rates by controlling endogenous memory growth, avoiding eviction-induced restarts, and keeping delay bounded. This framing also separates two regimes: in underloaded regimes, especially near the load boundary, the goal is to approximate the fluid upper bound; in overloaded regimes, the experiments examine how effectively each scheduler exploits available service capacity rather than losing work to eviction and restart. Accordingly, the paper now uses token-level \emph{effective throughput} for completed decode-token service net of evictions, reports request-level effective completion rate in the experiments, uses the fluid model to identify the stability region and the memory required to sustain it, and plots latency versus arrival rate in Section 6.
    \item \textbf{Rewrote the model and notation.} Section 2 now introduces prefill, decode, batching, KV-cache growth, the memory constraint, and the policy space narratively. The main notation is collected in the Notation Summary appendix, and the stage notation is separated from prompt lengths and output lengths. We also clarified that the memory constraint applies to the KV caches of all in-system requests retained on the GPU, not merely to requests in the current batch.
    \item \textbf{Clarified the time model and its scope.} We expanded the discussion around Equation (1), explaining how the linear memory-dependent iteration-time model fits the multi-stage serving model with endogenous memory growth. We added real-GPU validation in the Additional Experiments appendix, including direct A100 timing measurements and end-to-end GPU runs implemented with SGLang~0.5.7, an open-source inference infrastructure package~\citep{sglang-github}. The main policy comparisons remain Vidur simulations configured for A100 hardware; the real-GPU experiments provide direct timing calibration for the simulator and implementation evidence for the WAIT and Nested WAIT scheduling logic. We also explain why the model is well matched to prefill-decode disaggregated systems such as Splitwise and DistServe~\citep{patel2024splitwise,zhong2024distserve}, where separating prefill from decode makes aggregate KV-cache usage a natural state variable for iteration time. For richer calibrated service-time curves, the formal guarantees would require corresponding model-specific analysis, but the same scheduling perspective remains useful: admission and batching decisions must control the batch composition across prefill and decode stages because this composition drives both memory pressure and future iteration times.
    \item \textbf{Clarified the theoretical contribution.} We removed the misleading ``heavy traffic'' terminology from the narrative and now describe the guarantees as asymptotic guarantees. The main text now explains why the problem differs from classical formulations with exogenous service requirements: KV-cache memory grows endogenously with decode progress, evictions can restart partially completed work, and unknown output lengths are revealed only through the service trajectory. These features make the primitive state-action space substantially larger, because the scheduler must control both current batch composition and future memory growth. Without careful scheduling, evictions can occur frequently and waste effective throughput through repeated restart work. The WAIT and Nested WAIT thresholds reduce this high-dimensional state-action space to threshold and boundary queues that capture the memory-coupled dynamics.
    \item \textbf{Expanded and reorganized the experiments.} Section 6 now reports latency and effective-completion-rate curves over arrival-rate grids that include underloaded, near-overloaded, and overloaded regimes. The Additional Experiments appendix adds simulator-to-GPU validation on an A100, supplemental end-to-end real-GPU experiments, repeated-run variability near the transition from near-overloaded to overloaded, threshold-without-waiting comparisons, and a prefill-decode-disaggregated deployment study. For the real-data experiment, we compare the latency of different policies and use the effective-completion-rate curve to show how much offered load each policy can sustain. The real-GPU appendix further reports end-to-end scheduling performance on an A100 GPU using SGLang.
    \item \textbf{Updated positioning and scope.} The Introduction and related-work discussion now clarify that our contribution is a modeling and scheduling contribution for multi-stage LLM inference with endogenous memory growth, not a claim of a new heavy-traffic limit. We also added discussion of concurrent theoretical work on LLM inference scheduling and positioned our model as complementary: our focus is on memory-dependent iteration times, GPU KV-cache growth, and eviction-induced restarts under the partial information available to the scheduler.
    \item \textbf{Improved exposition and organization.} The comments on exposition led to a substantial rewrite of Sections 1--6 and the appendices. We removed ambiguous terminology, introduced notation more gradually, added more detailed explanations and instances to show the importance of the problem and where the main interests and difficulties arise, and reorganized the appendix proofs to improve clarity.
\end{itemize}

\begin{center}
\begin{tabularx}{0.96\textwidth}{@{}p{0.27\textwidth}X@{}}
\toprule
\textbf{Main concern} & \textbf{How the revised manuscript addresses it} \\
\midrule
Objective with exogenous arrivals & Clarifies the objective under exogenous arrivals with endogenous memory growth and a multi-stage decode mechanism: sustaining the offered load when feasible, avoiding eviction-induced restarts, and evaluating effective throughput, latency, and TTFT. \\
GPU KV-cache memory mechanism & Defines memory as the KV-cache usage of all admitted in-system requests retained on the GPU, distinguishes preemption from eviction, and makes restart costs explicit. \\
Linear time model & Explains the linear iteration-time model through aggregate KV-cache usage, reports A100 validation results, and discusses how the same state variable and scheduling insight can extend to prefill-decode-disaggregated deployments and related serving settings. \\
Theoretical contribution & Positions WAIT/Nested WAIT as fluid-guided threshold policies that control endogenous memory growth across prefill and decode stages, using threshold and boundary queues to reduce the high-dimensional state-action space under partial information. \\
Numerical experiments & Reports latency and effective-completion-rate curves over arrival-rate grids, adds long-decode and real-data experiments, repeated-run diagnostics near the transition from near-overloaded to overloaded, and supplemental real-GPU validation. \\
Writing and organization & Rewrites Sections 1--6, adds notation summary and examples, and reorganizes the appendices to improve proof clarity. \\
\bottomrule
\end{tabularx}
\end{center}

